% Shared by both templates: storage for the content fields, the private data,
% flags and icons. Compile with XeLaTeX (see build.mjs).
\usepackage{etoolbox}
\usepackage{pgfkeys}
\usepackage{graphicx}
\usepackage{tikz}
\usepackage{fontawesome5}
\usepackage{academicons}
\usepackage{enumitem}
\usepackage[hidelinks]{hyperref}

% --- Content fields --------------------------------------------------------
\newcommand\cvnombre[2]{\def\cvNombre{#1}\def\cvApellidos{#2}}
\newcommand\cvpuesto[1]{\def\cvPuesto{#1}}
\newcommand\cvenfoque[1]{\def\cvEnfoque{#1}}
\newcommand\cvresumen[1]{\def\cvResumen{#1}}
\newcommand\cvubicacion[1]{\def\cvUbicacion{#1}}
\newcommand\cvemail[1]{\def\cvEmail{#1}}
\newcommand\cvlinkedin[1]{\def\cvLinkedin{#1}}
\newcommand\cvgithub[1]{\def\cvGithub{#1}}
\newcommand\cvorcid[1]{\def\cvOrcid{#1}}
\newcommand\cvweb[1]{\def\cvWeb{#1}}
\newcommand\cvtelefono[1]{\def\cvTelefono{#1}}
\newcommand\cvnacimiento[1]{\def\cvNacimiento{#1}}

% Section titles: \cvtitulos{clave=Título, ...} and \cvt{clave}.
\pgfkeys{/cvt/.unknown/.code={\csdef{cvt@\pgfkeyscurrentname}{#1}}}
\newcommand\cvtitulos[1]{\pgfqkeys{/cvt}{#1}}
\newcommand\cvt[1]{\csuse{cvt@#1}}

% --- Private data ----------------------------------------------------------
% privado.tex (git-ignored) holds the phone number and birth date. The public
% build (\cvpublico) leaves them out even when the file exists.
\providetoggle{cvpublico}
\ifdefined\cvpublic\toggletrue{cvpublico}\fi
\IfFileExists{privado.tex}{\input{privado.tex}}{}
\iftoggle{cvpublico}{\let\cvTelefono\relax\let\cvNacimiento\relax}{}
\newcommand\siTelefono[1]{\ifdefined\cvTelefono\ifx\cvTelefono\relax\else#1\fi\fi}
\newcommand\siNacimiento[1]{\ifdefined\cvNacimiento\ifx\cvNacimiento\relax\else#1\fi\fi}

% --- Flags -----------------------------------------------------------------
\newcommand\bandera[1]{%
  \begin{tikzpicture}[x=1mm,y=1mm,baseline=0.6mm]
    \ifstrequal{#1}{es}{%
      \fill[red!80!black] (0,0) rectangle (9,6);
      \fill[yellow!85!orange] (0,1.5) rectangle (9,4.5);
      \fill[red!70!black] (2,2.4) rectangle (3.2,3.6);
    }{%
      \fill[blue!50!black] (0,0) rectangle (9,6);
      \begin{scope}
        \clip (0,0) rectangle (9,6);
        \draw[white,line width=1.2mm] (0,0) -- (9,6) (0,6) -- (9,0);
        \draw[red!80!black,line width=0.4mm] (0,0) -- (9,6) (0,6) -- (9,0);
      \end{scope}
      \fill[white] (3.6,0) rectangle (5.4,6) (0,2.1) rectangle (9,3.9);
      \fill[red!80!black] (3.95,0) rectangle (5.05,6) (0,2.45) rectangle (9,3.55);
    }
  \end{tikzpicture}}

\setlist{nosep}
